% Generated draft: requires mathematical and source review.
\documentclass[11pt]{article}
\usepackage[margin=25mm]{geometry}
\usepackage{fontspec}
\setmainfont{Latin Modern Roman}
\usepackage{amsmath,amssymb,mathtools,unicode-math}
\setmathfont{Latin Modern Math}
\usepackage{enumitem,needspace,xurl,hyperref}
\setlength{\parindent}{0pt}
\setlength{\parskip}{5pt}
\setlength{\emergencystretch}{3em}
\newcommand{\N}{\mathbb{N}}
\newcommand{\Z}{\mathbb{Z}}
\newcommand{\Q}{\mathbb{Q}}
\newcommand{\R}{\mathbb{R}}
\newcommand{\C}{\mathbb{C}}
\newcommand{\F}{\mathbb{F}}
\newcommand{\A}{\mathbb{A}}
\newcommand{\PP}{\mathbb{P}}
\newcommand{\E}{\mathbb{E}}
\newcommand{\eps}{\varepsilon}
\newcommand{\dd}{\,\mathrm{d}}
\begin{document}
\section*{UnsolvedMath ID: 30006434}
\textbf{Source problem code:} catalogue entry
\subsection{Definitions and mathematical statement}
Let \(\mathcal A_d\) denote the intended class of \(d\)-dimensional algebras in the subalgebra-zeta setting, and write \(Z_A^{\mathrm{red}}(T)\) for the reduced subalgebra zeta function attached to \(A\). The supplied record does not specify this class, the ground ring of \(A\), the meaning of its dimension, or the construction of the reduced zeta function. These are unresolved parameters of the statement.

For a field \(k\) and an \(\N_0\)-graded \(k\)-algebra \(R=\bigoplus_{n\geq 0}R_n\) with \(\dim_k R_n<\infty\), its Hilbert series is
\[
\operatorname{Hilb}_{R,k}(T)=\sum_{n\geq 0}(\dim_k R_n)T^n.
\]
Under the usual commutative Noetherian interpretation, \(R\) is Gorenstein if every localization \(R_{\mathfrak p}\) at a prime ideal has finite injective dimension as a module over itself; its dimension means Krull dimension. The coefficient field and the precise algebraic, finiteness, and connectedness conventions intended for \(R\) are not supplied, so this interpretation requires confirmation.

\textbf{Extracted conjectural target.} For each dimension \(d\) allowed by the intended setting and every \(A\in\mathcal A_d\), construct a natural \(\N_0\)-graded Gorenstein algebra \(R_A\) of dimension \(d\) such that
\[
\operatorname{Hilb}_{R_A,k}(T)=Z_A^{\mathrm{red}}(T).
\]
Here the equality means equality of formal power series at \(T=0\), once the relevant series and coefficient field have been specified. The word ``natural'' is retained from the extracted statement; no required functoriality or canonical construction is defined in the record. Consequently, this is a statement template with explicit unresolved qualifications, not a fully specified universal conjecture.

\textbf{Supplied qualifications and context.} The excerpt attributes the Hilbert-series conjecture to Voll and says that it circulated informally while published formulations were more cautious. The excerpt does not reproduce the full conjecture or its hypotheses, so it does not by itself establish the scope of the extracted target above.

The excerpt reports that, when \(A\) admits a particular type of basis, Evseev's result, cited there as reference [6, Proposition 4.1], produces a \(d\)-dimensional cone \(C\subseteq\R_{\geq 0}^{d}\) for which \(Z_A^{\mathrm{red}}(T)\) is the coarse Hilbert series of
\[
\Q[C\cap\Z^d].
\]
This monoid algebra has basis \(\{x^u:u\in C\cap\Z^d\}\) and multiplication \(x^ux^v=x^{u+v}\). The excerpt further reports that Evseev's explicit cone description, together with results cited as [11, Chapter 1], makes this algebra Gorenstein. The basis condition, cone description, and coarse grading are absent from the supplied record; this restricted result must not be extended to all \(A\).

The excerpt also states that the Hilbert-series conjecture would imply a pole of order \(d\) at \(T=1\), and hence a finite nonzero value
\[
m_{\mathrm{red}}(A)=\lim_{T\to 1}(1-T)^d Z_A^{\mathrm{red}}(T).
\]
It separately relates a topological degree conjecture to finiteness and nonvanishing of a quantity \(m_{\mathrm{top}}(A)\). The displayed definition of that quantity is corrupted in the supplied extraction and cannot be specified reliably here.
\subsection{Short English statement}
The extracted conjecture asks for a natural graded Gorenstein algebra of dimension \(d\) whose Hilbert series equals the reduced subalgebra zeta function of each eligible \(d\)-dimensional algebra \(A\). The eligible algebras, reduced zeta construction, coefficient field, and meaning of naturality remain unspecified. The supplied excerpt reports a restricted cone construction and attributes the broader conjecture to Voll; a complete formulation requires human review.
\subsection{Sources}
\begin{itemize}
\item Supplied local UnsolvedMath registry record, permanent ID \(30006434\), source problem code \texttt{catalogue entry}; registry path \texttt{data/unsolvedmath/problems.json}. This draft uses the supplied record without independently inspecting the registry or external sources.
\item Supplied report citation: \emph{Cohomology Theories for Automorphic Forms and Enumerative Algebra} (2026), \emph{Oberwolfach Reports}, pp.~2065--2124, DOI \texttt{10.4171/owr/2025/38}. \url{https://doi.org/10.4171/owr/2025/38}. The record separately gives proposed year 2025; the relationship between these dates is not explained.
\item Supplied literature source, ``Oberwolfach Report 38'': \url{https://ems.press/content/serial-article-files/52251}.
\item Supplied literature source, ``Rossmann--Voll survey slides'': \url{https://torossmann.github.io/files/mfo25b.pdf}.
\item Attribution in the supplied excerpt: Voll's Hilbert-series conjecture; Evseev, reference [6, Proposition 4.1]; combinatorial commutative algebra results, reference [11, Chapter 1]. Bibliographic details for [6] and [11] are not supplied, and no individual proposer is identified in the registry.
\item Literature assessment dated exactly \texttt{2026-08-22}: the supplied curation labels the problem open, describes the general Hilbert-series conjecture as unresolved, and describes Evseev's restricted construction as partial progress. This is supplied dated curation, not an independently verified literature assessment. The registry also supplies the label \texttt{corrected\_verified} and a claim that the official report was checked; those are provenance claims of the record, not verification performed for this draft.
\end{itemize}
\end{document}
